\documentclass[10pt]{article}
\usepackage[margin=0.68in]{geometry}
\usepackage{amsmath,amssymb,amsthm,graphicx,microtype,xcolor,hyperref,placeins,booktabs,multicol}
\hypersetup{colorlinks=true,linkcolor=blue!45!black,citecolor=blue!45!black,urlcolor=blue!55!black}
\setlength{\parindent}{0pt}
\setlength{\parskip}{3.5pt}
\setlength{\floatsep}{8pt plus 2pt minus 2pt}
\setlength{\textfloatsep}{10pt plus 2pt minus 2pt}
\setlength{\intextsep}{8pt plus 2pt minus 2pt}
\renewcommand{\topfraction}{0.90}
\renewcommand{\textfraction}{0.08}
\renewcommand{\floatpagefraction}{0.70}
\newcommand{\E}{\mathbb E}
\newcommand{\Prob}{\mathbb P}
\newcommand{\dd}{\,\mathrm d}
\newcommand{\valfig}{C:/Users/anees/OneDrive/Documents/GitHub/local_time_traps/IRIS paper/IRIS synopsis/figures/monte_carlo_convergence.png}
\newcommand{\appfig}{C:/Users/anees/OneDrive/Documents/GitHub/local_time_traps/IRIS paper/IRIS synopsis/figures/fig03_state_of_practice_vs_point_trap.pdf}
\newtheorem{theorem}{Theorem}
\newtheorem{conjecture}{Conjecture}

\title{\textbf{Point-Trap Design in Diffusive Media: Minimax Capture, Graph Bounds, and Spectral Placement}}
\author{Aneesh Arnav Chikkala}
\date{September 28, 2026}

\begin{document}
\maketitle

\begin{abstract}
	In many transport settings, a diffusing particle is removed only when it encounters one of a limited number of reactive sites. The placement of these sites is therefore an optimization problem: with a fixed total amount of reactivity, the designer chooses their locations and strengths to control the longest capture time over all starting positions.
	
My study examines reflected diffusion in a heterogeneous one-dimensional medium with finitely many partially absorbing point traps. My main result reduces the continuum placement problem to a finite source-partition problem. A balancing principle equalizes resistance-weighted rises in each cell, and a monotone budget equation determines the strengths. The source share $p_i$, barrier $R_i$, and strength $\kappa_i$ satisfy $p_iR=\kappa_i(T-R_i)$, so source share is proportional to strength times remaining barrier capacity.
	
	I also derive a mass--resistance bound for metric graphs and a spectral result for long-time survival with prescribed strengths. In a saturated-sediment benchmark, the optimized placement lowers the modeled worst-start lifetime by $4.246\%$ relative to equal physical-depth placement. The interval theorem is formally checked in Lean under its stated assumptions, while the framework provides a testbed for related open problems, including narrow-capture and spectral-partition questions connected with the Hot Spots Conjecture. Together, these results provide a route from heterogeneous transport coefficients to optimized trap placement, with the sediment benchmark illustrating a possible application.
\end{abstract}
\section{Introduction and optimization problem}

Let $X_t$ be a reflected diffusion on $[0,L]$ with transport coefficient $a(x)>0$ and stored-source density $\rho(x)>0$. A design $\mathcal D$ consists of ordered point traps $0<x_1<\cdots<x_N<L$ with strengths $\kappa_j\geq0$ and fixed total capacity $\sum_j\kappa_j=K$. The traps are partially reactive: capture occurs through repeated local-time contact rather than instantaneous absorption.

Let $\tau_{\mathcal D}$ be the first-capture time and $u_{\mathcal D}(x)=\E_x[\tau_{\mathcal D}]$. The design problem is the fixed-budget minimax problem
\[
M_N^*(K)=\inf_{\mathcal D:\,\sum_j\kappa_j=K}\max_{x\in[0,L]}u_{\mathcal D}(x).
\]
It asks for the placement and allocation that protect the least favorable starting point.

The interval reduction, graph equality conditions, and spectral comparison address different observables. The discounted objective has an additional zeroth-order term and is treated separately.

\section{Innovation and contributions to open problems}

\subsection{Constructive reduction}

For any admissible layout, the trap sites determine source shares. Balancing the two rises in each cell cannot increase the worst-start lifetime. The remaining strength allocation is then obtained from one monotone budget root. Conversely, every admissible partition reconstructs a physical layout with the reduced value. In compact notation,
\[
\inf_{\mathcal D}M(\mathcal D)=\min_{q\in\Delta_{N-1}}T(q),
\]
where $T(q)$ is the common balanced barrier and $\Delta_{N-1}$ is the compactified source-partition simplex.

\subsection{Open-problem contributions}

The interval result retains heterogeneous coefficients through the source-coordinate inverse and resistance density. On metric graphs, the lower bound separates three losses: reactive allocation, edge-profile curvature, and mismatch between source mass and resistance. The spectral calculation places prescribed strengths by Rayleigh comparison and a one-jump phase inequality. Finite-width reactions and the discounted equal-rise inequality remain open under advection or heterogeneous coefficients.

\subsection{Proposed conjecture: discounted capture}

Let $T_R$ be an independent exponential deadline with rate $s$. Since
\[
\Prob_x(\tau<T_R)=\E_x[e^{-s\tau}],
\]
the discounted objective is the probability of capture before an unpredictable disturbance or reporting deadline. For a homogeneous reflecting interval with symmetric equal-strength traps, the observed formula suggests
\[
\sup_{\mathcal D}\inf_x\E_x[e^{-s\tau_{\mathcal D}}]
\stackrel{?}{=}
\left[\cosh\eta+\frac{2N\sqrt{Ds}}{K}\sinh\eta\right]^{-1},
\qquad
\eta=\frac{L}{2}\sqrt{\frac{s}{D}}.
\]
\begin{conjecture}
For every $s>0$, the symmetric equal-strength midpoint design maximizes worst-start discounted capture over all placements and positive strength allocations with fixed count and total strength on a homogeneous reflecting interval.
\end{conjecture}

The zero-discount limit recovers the mean-time optimizer, while large $s$ emphasizes immediate contact. These observations motivate the conjecture. Its proof requires an equal-rise inequality for the equation with the additional zeroth-order term; heterogeneous coefficients, advection, and finite-width reactions are outside its scope.

\section{Methodology and analytical framework}

\subsection{Local-time model and backward equation}

Let $a(x)>0$ be the transport coefficient and $\rho(x)>0$ the stored source density. For trap sites $0<x_1<\cdots<x_N<L$ and strengths $\kappa_j\geq0$ satisfying $\sum_j\kappa_j=K$, the backward equation for the mean lifetime is
\[
\frac{1}{\rho(x)}\frac{\mathrm d}{\mathrm dx}\left(a(x)u'(x)\right)=-1
\quad\text{away from traps},
\]
with reflecting boundary conditions $a(0)u'(0)=a(L)u'(L)=0$ and interface jumps
\[
\left[a u'\right]_{x_j^-}^{x_j^+}=\kappa_j u(x_j).
\]
The stochastic representation is local-time killing,
\[
\Prob_x(\tau_{\mathcal D}>t)=
\E_x\!\left[\exp\left(-\sum_{j=1}^N\kappa_jL_t^{x_j}\right)\right],
\]
with local-time normalization fixed by the jump condition. A particle may revisit a weak trap many times before capture; this is why the point interaction is not simply a zero-width version of instantaneous absorption.

\subsection{Source coordinates and equal-rise cells}

Set
\[
q(x)=\frac{1}{R}\int_0^x\rho(s)\dd s,
\qquad R=\int_0^L\rho(s)\dd s,
\qquad
\omega(q)=\frac{R^2}{a(x(q))\rho(x(q))}.
\]
The transformed coordinate $q\in[0,1]$ records source mass; the inverse $\mathrm dx/\mathrm dq=R/\rho(x(q))$ returns physical distances after optimization. For a source cell $[\alpha,\beta]$ and trial trap location $c$, define the one-sided rises
\[
A(c)=\int_\alpha^c(c-s)\omega(s)\dd s,
\qquad
B(c)=\int_c^\beta(s-c)\omega(s)\dd s.
\]
Because $A'(c)>0$ and $B'(c)<0$, $A(c)-B(c)$ is strictly increasing. There is one equal-rise point $c_i$ and one local minimax barrier $R_i=A(c_i)=B(c_i)$. This is the analytical reason that an unbalanced cell cannot be optimal: moving the trap toward the larger rise lowers the active maximum.

Write $\widehat K:=K/R$ for the normalized budget. If $p_i$ is the source share of cell $i$, equalization at common barrier $T$ gives
\[
p_iR=\kappa_i(T-R_i),
\qquad
\kappa_i(T)=\frac{Rp_i}{T-R_i},
\qquad T>\max_iR_i.
\]
The total-strength constraint is the scalar equation
\[
\Phi(T):=\sum_{i=1}^N\frac{p_i}{T-R_i}=\widehat K,
\qquad\text{equivalently}\qquad\sum_i\kappa_i(T)=K.
\]
Since $\Phi'(T)=-\sum_i p_i/(T-R_i)^2<0$, the root is unique, bracketable, and suitable for certified bisection. The reduced objective is continuous on the compactified partition simplex, so the outer minimum is attained. For any physical design, $M-u(x_i)\ge R_i$ and $Rp_i=\kappa_i u(x_i)$, hence $\Phi(M)\le K/R$; monotonicity gives $T\le M$. The reconstructed equal-rise design attains equality, proving the global reduction.

The root also supplies sensitivity information. With $w_i=p_i/(T-R_i)^2$ and $S=\sum_iw_i$,
\[
\frac{\partial T}{\partial K}=-\frac{1}{RS},
\qquad
\frac{\partial T}{\partial R_j}=\frac{w_j}{S},
\qquad
\delta T=\frac{\sum_iw_i\delta R_i-\delta K/R}{S}+o(\lVert(\delta R,\delta K)\rVert).
\]
Thus added strength lowers the barrier, while a geometrically difficult cell raises it in proportion to its budget sensitivity.

\subsection{Graph and spectral companions}

On a compact metric graph, edge $e$ has length $\ell_e$, conductance $a_e$, and source density $\rho_e$. The weighted operator and vertex condition are
\[
Hf=-\frac{1}{\rho_e}(a_ef')',
\qquad
\sum_{e\sim v}a_e\partial_\nu f(v)=J_vf(v).
\]
Integrating the torsion equation $Hu=1$ over edges gives conservation of source mass. Comparing the torsion profile with its edgewise linear profile yields
\[
M\geq\frac{V}{K}+\frac{V}{8C},
\qquad
\Delta_{\mathrm{graph}}:=M-\frac{V}{K}-\frac{V}{8C}
=\Delta_{\mathrm{react}}+\Delta_{\mathrm{edge}}+\Delta_{\mathrm{mass}}\geq0.
\]
The three terms identify, respectively, reaction allocation loss, curvature loss, and mass--resistance imbalance. Equality requires all active cells to balance simultaneously, which explains why passive leaves or mismatched parallel edges create a strict gap.

For the spectral companion, the killed operator has quadratic form
\[
\mathcal Q(v)=\int a|v'|^2\dd x+\sum_j\kappa_j|v(x_j)|^2,
\qquad
\lambda_1=\inf_{v\neq0}\frac{\mathcal Q(v)}{\int\rho|v|^2\dd x}.
\]
The Rayleigh principle and a phase-jump comparison control the principal decay rate, a separate observable from $\max_xu(x)$.

\section{Results, validation, and real-world application}

\subsection{Analytical results}

The interval theorem proves existence and gives the exact minimax value as the minimum common barrier over source partitions. In a homogeneous interval with diffusivity $D$, length $L$, $N$ traps, and total strength $K$, the unique optimizer has equal strengths at the midpoints of equal subintervals, with
\[
\|u\|_\infty=\frac{L}{K}+\frac{L^2}{8DN^2}.
\]
The first term is the reaction-budget cost; the second is the diffusion-and-spacing cost. Adding traps reduces only the second term. In heterogeneous media, the source-coordinate construction replaces equal physical spacing by equalized source shares and resistance-aware barriers.

The graph and spectral results give separate extensions: graph equality conditions and prescribed-strength placement. They also identify failure modes for a naive equal-spacing rule, including unequal source density, passive branches, and finite-width reactions [3--5].

\subsection{Real-world application benchmark}

We now propose an application in a dimensionless saturated porous slab. Dissolved pollutant diffuses through porewater, and two narrow activated-carbon bands act as localized sinks. The reproducible setup is $L=1$, $a(x)=1/(1+0.85x)$, $\rho(x)=1+0.55x$, $K=2$, two equal-width bands of width $0.055$, and equal strengths.

For a first-pass remediation design, equal physical-depth placement is the baseline. It ignores heterogeneity in source mass and transport resistance. The proposed rule uses the measured medium profile to select a source-balanced layout before an independent numerical comparison.

We keep total strength and band width fixed. The comparison uses equal-depth placement, the theorem-guided layout, and an independent finite-volume calculation. The recorded sites are $(.250,.750)$, $(.306,.792)$, and $(.296,.789)$, respectively; the source-coordinate comparator is $(.384,.711)$.

For two cells, let $q_1$ be the source share of the left cell, so $p_1=q_1$ and $p_2=1-q_1$. If $R_i(q_1)$ is the equal-rise barrier in cell $i$, the guided barrier $T$ is the unique root above both barriers of
\[
\frac{q_1}{T-R_1(q_1)}
+\frac{1-q_1}{T-R_2(q_1)}=K/R.
\]
Equivalently,
\[
\frac{K}{R}(T-R_1)(T-R_2)=q_1(T-R_2)+(1-q_1)(T-R_1).
\]
Equal physical depth fixes $x_1=L/2$ and therefore fixes $q_1=q(L/2)$, which need not be $1/2$ when $\rho$ is heterogeneous. The theorem-guided design minimizes the resulting $T(q_1,K)$ over $q_1$ and reconstructs $x_1=q^{-1}(q_1)$. The gain comes from matching source mass and resistance.

For the declared heterogeneous coefficient profile, the equal-depth baseline has worst-start lifetime $M_{\mathrm{base}}=0.722868$. The analytic design gives $M_{\mathrm{guided}}=0.692174$, a $4.246\%$ reduction. The finite-volume reference gives $0.687172$, a $4.938\%$ reduction, leaving a gap of $0.005002$. The midpoint early-time capture diagnostic changes from $0.469327$ to $0.470173$.

The two-cell root also gives a simple capacity check. With $q_1$ fixed and $w_i=p_i/(T-R_i)^2$,
\[
\frac{\partial T}{\partial K}=-\frac{1}{R(w_1+w_2)}.
\]
Thus extra total strength lowers the common barrier, while the source-share optimization changes the geometric terms $R_1(q_1)$ and $R_2(q_1)$. In the weak-reaction limit,
\[
T=\frac{R}{K}+q_1R_1+(1-q_1)R_2+O(K/R),
\qquad K\downarrow0,
\]
so the first correction is the source-weighted diffusion barrier.

At fixed $K$, relocation changes the geometric term. When $\rho$ varies, equal physical depth does not generally give equal source shares, so the source-coordinate rule changes the diffusion contribution while preserving total strength.

\begin{figure}[!ht]
\centering
\includegraphics[width=.80\linewidth]{\appfig}
\caption{Application benchmark. Equal physical-depth placement is compared with the theorem-guided point-trap design and an independent finite-volume reference.}
\label{fig:application}
\end{figure}
\FloatBarrier

The reduced model omits advection, nonlinear sorption, concentration-dependent uptake, dispersion, finite-rate bands, uncertainty in $a$ and $\rho$, and spatially variable access to amendment. The next experimental test is to calibrate these mechanisms and compare mass removed, threshold time, amendment cost, and worst-start lifetime.

\subsection{Implications, open problems, and translation}

For a new medium, the screening calculation needs the measured $q(x)$, $\omega(q)$, and total capacity $K$; it returns a physical placement and strength allocation for a detailed transport model.

The interval and graph results are proved. The discounted-capture statement remains open, with the equal-rise inequality as the next proof target.

The Lean audit checks the proof term under the recorded definitions, assumptions, imports, and environment. The certificate therefore applies to the encoded theorem and its assumptions; application mechanisms still require calibration.

After calibration of $a(x)$ and $\rho(x)$, the same construction can be compared with a finite-width transport model using mass removed, threshold time, and amendment cost as additional outputs.

\subsection{Validation protocol and scope}

The validation record combines grid refinement, residuals, an independent discretization, the graph identity, and the Lean kernel check. These records address numerical stability, identity checks, cross-discretization agreement, and the formal proof environment.

\begin{samepage}
For a positive discrete operator $A_h$, source vector $b_h$, approximate solution $\widetilde u_h$, and residual $r_h=A_h\widetilde u_h-b_h$, the exact discrete solution satisfies
\[
u_h-\widetilde u_h=-A_h^{-1}r_h,
\qquad
\|u_h-\widetilde u_h\|_\infty
\leq\|A_h^{-1}\|_\infty\|r_h\|_\infty.
\]
The inverse-operator factor controls the residual-to-error estimate. Continuum accuracy requires refinement or an independent discretization.
\end{samepage}

The validation suite reports a maximum linear-system residual of $1.3378\times10^{-14}$ over $5{,}000$ cases, a maximum graph-deficit residual of $3.6082\times10^{-16}$ over $2{,}400$ cases, and an independent finite-element relative error of $1.4275\times10^{-8}$. Full records are in \path{verification/README.md}, \path{verification/results/}, and \path{verification/generalized/results/}.

\pagebreak
\begin{samepage}
The formal route is in module \texttt{TimeTraps.Interval}. It includes \texttt{global\_minimum} in \texttt{IntervalMedium}; \newline\texttt{optimal\_iff\_reconstruction} in \texttt{IntervalDesign}. The proof map is \path{formalization/review/INTERVAL.md}; the certificate is \path{formalization/audit/current-verification/result.json}. The pinned environment is Lean 4.33.1 with Mathlib revision \texttt{0df444a3}; the audit records $2{,}282$ research modules, $10{,}543$ named theorems, and $17{,}355$ audited declarations.
\end{samepage}

\begin{figure}[!ht]
\centering
\includegraphics[width=.55\linewidth]{\valfig}
\caption{Convergence diagnostic with horizon $2$ and $6{,}000$ paths. As $\Delta t$ decreases from $0.002$ to $0.00025$, the estimates rise from $0.25310$ to $0.25557$ while remaining below the finite-horizon target $0.2591433$ and full mean $0.2592593$. The below-target trend records timestep, finite-band, and omitted-tail bias.}
\label{fig:validation}
\end{figure}
\FloatBarrier

The theorem supplies the placement rule; the validation record documents its implementation. Calibration and the discounted-capture conjecture remain open.

\section{Acknowledgement and references}

The author thanks the open-source scientific-computing and formalization communities. The environmental interpretation is a mathematical benchmark, and the conjecture and application protocol are offered as testable next steps.

\begin{multicols}{2}
\footnotesize
\begin{enumerate}
\setlength{\itemsep}{2pt}
\item P. C. Bressloff, \emph{J. Phys. A} 55 (2022), 205001.
\item U. Ghosh et al., \emph{Environ. Sci. Technol.} 45 (2011), 1163--1168.
\item G. Buttazzo et al., \emph{Ann. Mat. Pura Appl.} 185 (2006), 207--221.
\item D. Mugnolo and M. Pl{"u}mer, \emph{Calc. Var.} 62 (2023), 27.
\item A. A. Chikkala, ``Optimal Design of Partially Reactive Point Traps,'' preprint (2026).
\item A. A. Chikkala, \emph{Point-Trap Design in Diffusive Media}, research archive containing the journal paper, Lean code, and handwritten mathematics, GitHub (2026). \url{https://github.com/aneesharnavch/Point-Trap-Design-in-Diffusive-Media}
\end{enumerate}
\end{multicols}

\end{document}
